\documentclass[10pt,a4paper]{amsart}
\usepackage[margin=19mm]{geometry}
\usepackage{amsmath,amssymb,mathtools}
\usepackage[scr=rsfs,cal=euler]{mathalfa}
\usepackage{xfrac}
\usepackage[unicode,hidelinks,bookmarks=false]{hyperref}
\newcommand{\ie}{i.e.\ }
\newcommand{\cf}{cf.\ }
\newcommand{\resp}{respectively\ }
\newcommand{\Subsection}{\subsection}
\providecommand{\nobreakdash}{\nobreak\hbox{-}}
\setlength{\emergencystretch}{2em}
\multlinegap0pt
\usepackage{bm}
\usepackage{bbm}
\usepackage{dsfont}
\usepackage{xspace}
\usepackage{cancel}
\usepackage{mathtools}
\usepackage{amsthm}
\makeatletter
\@ifundefined{theorem}{\newtheorem{theorem}{Theorem}}{}
\makeatother
\title[Novel Perturbed b-Metric and Perturbed Extended b-Metric Spa]{Novel Perturbed b-Metric and Perturbed Extended b-Metric Spaces with Banach-Type Fixed Point Theorem}
\author{Abdelhamid Moussaoui}
\date{Source: arXiv:2512.21341v1 (CC BY 4.0)}
\begin{document}
\maketitle

\noindent\emph{Source and licence.} Novel Perturbed b-Metric and Perturbed Extended b-Metric Spaces with Banach-Type Fixed Point Theorem. Version arXiv:2512.21341v1 (CC BY 4.0), distributed under the Creative Commons Attribution 4.0 International licence, as recorded in the arXiv licence metadata for this exact version and shown on the abstract page https://arxiv.org/abs/2512.21341v1. Full rights notice: licenses/arxiv-2512.21341-CC-BY-4.0.txt.\par\smallskip

\noindent\textit{Editorial scope.} Counted unit: the Theorem whose source label is \texttt{3.1} in the article (source0.tex lines 626--642, source label \texttt{3.1}), delivered together with the article's own complete original proof of it (source0.tex lines 644--814, one \texttt{proof} environment of 171 lines). No mathematics is written, repaired or re-derived: statement and proof are the article's own text line for line. Required context retained uncounted: none. Every definition, display and label of those lines is retained unchanged. Omitted from this package and not redistributed: 1--625, 643--643, 815--1065. Nothing inside a retained excerpt is omitted or altered: every retained body is delivered exactly as the article's lines are written, so no omission receipt of any kind accompanies this package. Citations: the retained excerpts carry no citation command at all, so no bibliography entry of the article is transcribed and no reference is redistributed. Reference adaptation: none was needed and none was made; every reference inside the delivered unit resolves inside it, so no unresolved reference or citation is produced. Printed slips kept literal: no printed word, symbol, hypothesis or number of the counted statement or of its proof was altered. Layout and class: the article's own class is \texttt{sn-jnl} with packages including graphicx, multirow, amsmath, amssymb, amsfonts, amsthm, mathrsfs, appendix, xcolor, textcomp, manyfoot, booktabs, algorithm, algorithmicx; the wrapper is the lane's pinned \texttt{amsart} editorial wrapper at 10pt on A4, which restates the theorem-like environments the retained text needs and adds the article's own macro definitions that the retained text needs (none), with extra wrapper packages (bm, bbm, dsfont, xspace, cancel, mathtools). The delivered numbering is the unit's own. Assets: no retained span contains a figure environment, a graphics-inclusion command, a PDF-page inclusion, a table or a TikZ picture; no image file is copied into this package, the package declares no asset dependency and no image file is redistributed with it. Licence: version arXiv:2512.21341v1 of this article is distributed under the Creative Commons Attribution 4.0 International licence, as recorded in the arXiv licence metadata for this exact version and shown on the abstract page https://arxiv.org/abs/2512.21341v1. A full rights notice is delivered at licenses/arxiv-2512.21341-CC-BY-4.0.txt. Characters: every retained excerpt is pure ASCII, so the delivered engine, XeLaTeX with its default Latin Modern text font, reports no missing character.
\begin{theorem}\label{3.1} 
Let $(X,\mathcal{D}_{\zeta},\hbar)$ be a complete perturbed extended $b$-metric space. Assume that the associated exact extended b-metric is a continuous function. Let $T : X \to X$ be a mapping satisfying
\begin{equation}
\mathcal{D}_{\zeta}(T v,T \omega) 
\le c.\, \mathcal{D}_{\zeta}(v,\omega),
\qquad \text{for all } v,\omega \in X,
\end{equation}
for some constant $c \in [0,1)$. 

Moreover, suppose that for each $v_{0} \in X$, the sequence 
$v_{n} = T^{n}v_{0}$, $n = 1,2,\dots$, satisfies
\[
\lim_{n,m \to +\infty} \zeta(v_{n},v_{m}) < \frac{1}{c}.
\]
Then $T$ has a unique fixed point $\vartheta \in X$. Furthermore, for every 
$v \in X$, the sequence $(T^{n}v)$ converges to $\vartheta$ with respect to $\mathcal{D}_{\zeta}$.
\end{theorem}

\begin{proof}
Let $v_{0}\in X$ be arbitrary and construct the iterative sequence 
$\{v_{n}\}$ by
\[
v_{0}, \qquad v_{1}=T v_{0}, \qquad v_{2}=T v_{1}=T^{2}v_{0},\quad \dots,\quad 
v_{n}=T^{n}v_{0},\ \dots
\]

From the contractive condition
\[
\mathcal{D}_{\zeta}(T v,T \omega) \le c\,\mathcal{D}_{\zeta}(v,\omega)
\qquad \text{for all } v,\omega\in X,
\]
we deduce, for $n=0$,
\[
\mathcal{D}_{\zeta}(v_{1},v_{2})
= \mathcal{D}_{\zeta}(T v_{0},T v_{1})
\le c\,\mathcal{D}_{\zeta}(v_{0},v_{1}).
\]

Assume now that for some $n\ge0$ we have
\[
\mathcal{D}_{\zeta}(v_{n},v_{n+1})
\le c^{n}\,\mathcal{D}_{\zeta}(v_{0},v_{1}).
\]
Using the contractive condition with $v=v_{n}$ and $\omega=v_{n+1}$, we obtain
\[
\mathcal{D}_{\zeta}(v_{n+1},v_{n+2})
= \mathcal{D}_{\zeta}(T v_{n},T v_{n+1})
\le c\, \mathcal{D}_{\zeta}(v_{n},v_{n+1})
\le c^{\,n+1}\,\mathcal{D}_{\zeta}(v_{0},v_{1}).
\]

Thus, by induction on $n$, we conclude that
\begin{equation}\label{Dz-iter}
\mathcal{D}_{\zeta}(v_{n},v_{n+1})
\le c^{\,n}.\Im ~~\text{ for all } n\in\mathbb{N},
\end{equation}
where $\Im=\mathcal{D}_{\zeta}(v_{0},v_{1}).$


Let \(
d_{\zeta}(u,w)=\mathcal{D}_{\zeta}(u,w)-\hbar(u,w)
\)
be the exact extended $b$-metric associated with $(X,\mathcal{D}_{\zeta},\hbar)$.  By the definition of $d_{\zeta}$, we may rewrite \eqref{Dz-iter} as
\[
d_{\zeta}(v_n,v_{n+1}) + \hbar(v_n,v_{n+1})
= \mathcal{D}_{\zeta}(v_n,v_{n+1})
\le c^{\,n}\,\mathcal{D}_{\zeta}(v_0,v_1), \qquad n\in\mathbb{N}.
\]

Since
\[
d_{\zeta}(v_n,v_{n+1})
\le d_{\zeta}(v_n,v_{n+1})+\hbar(v_n,v_{n+1})
= \mathcal{D}_{\zeta}(v_n,v_{n+1}),
\]
we deduce from \eqref{Dz-iter} that
\begin{equation}\label{dz-final}
d_{\zeta}(v_n,v_{n+1})
\le c^{\,n}.\Im,
\qquad n\in\mathbb{N}.
\end{equation}
%A standard argument applied to \eqref{dz-final} shows that the sequence 
%$\{v_n\}$ is Cauchy in the exact extended $b$-metric space $(X,d_{\zeta})$.  
%Since $(X,\mathcal{D}_{\zeta},\hbar)$ is complete (i.e., $(X,d_{\zeta})$ is complete), 
%there exists a point $\vartheta\in X$ such that
%\[
%\lim_{n\to+\infty} d_{\zeta}(v_n,\vartheta)=0.
%\]
%Thus the sequence $\{v_n\}$ converges to $\vartheta$ in the metric space $(X,d_{\zeta})$.
Using successively the extended $b$-metric inequality, for $m>n$ we have
\[
\begin{aligned}
d_{\zeta}(v_n , v_m)
&\le \zeta(v_n , v_m)\, d_{\zeta}(v_n , v_{n+1})
   + \zeta(v_n , v_m)\zeta(v_{n+1},v_m)\, d_{\zeta}(v_{n+1},v_{n+2})  \\
&\quad +\; \zeta(v_n , v_m)\zeta(v_{n+1},v_m)\zeta(v_{n+2},v_m)\,
          d_{\zeta}(v_{n+2},v_{n+3}) + \cdots \\
&\quad +\; \zeta(v_n , v_m)\zeta(v_{n+1},v_m)\cdots \zeta(v_{m-1},v_m)\,
          d_{\zeta}(v_{m-1},v_m).
\end{aligned}
\]

From the estimate
\[
d_{\zeta}(v_j , v_{j+1}) \le c^{\,j}\,\Im, \qquad j\in\mathbb{N},
\]
we deduce
\[
\begin{aligned}
d_{\zeta}(v_n , v_m)
&\le \Im \Big[
c^{n}\prod_{i=1}^{n}\zeta(v_i,v_m)
+ c^{\,n+1}\prod_{i=1}^{n+1}\zeta(v_i,v_m)
+ \cdots  \\
&\qquad\qquad 
+ c^{\,m-1}\prod_{i=1}^{m-1}\zeta(v_i,v_m)
\Big].
\end{aligned}
\]

Since 
\[
\lim_{n,m\to+\infty} \zeta(v_{n+1},v_m)\,c < 1,
\]
the series
\[
\sum_{j=1}^{+\infty} c^{\,j}\prod_{i=1}^{j}\zeta(v_i,v_m)
\]
converges by the ratio test for each fixed $m\in\mathbb{N}$.  
Define
\[
\daleth=\sum_{j=1}^{+\infty} c^{\,j}\prod_{i=1}^{j}\zeta(v_i,v_m),
\qquad
\daleth_n=\sum_{j=1}^{n} c^{\,j}\prod_{i=1}^{j}\zeta(v_i,v_m).
\]

Thus, for $m>n$, the above inequality gives
\[
d_{\zeta}(v_n , v_m)
\le \Im\,\big(\daleth_{m-1}-\daleth_n\big).
\]

Letting $n\to+\infty$, we conclude that 
\[
d_{\zeta}(v_n,v_m)\longrightarrow 0,
\]
hence the sequence $\{v_n\}$ is Cauchy in the exact extended $b$-metric space $(X,d_{\zeta})$. In particular, by definition of the perturbed structure, this means that $\{v_n\}$ is a perturbed Cauchy sequence in $(X,\mathcal{D}_{\zeta},\hbar)$. Since $(X,\mathcal{D}_{\zeta},\hbar)$ is complete in the perturbed sense, and 
$\{v_n\}$ has been shown to be Cauchy in $(X,d_{\zeta})$, there exists a point 
$\vartheta \in X$ such that
\begin{equation}\label{limit-v}
\lim_{n \to +\infty} d_{\zeta}(v_n,\vartheta)=0.
\end{equation}

\medskip
We now verify that $\vartheta$ is a fixed point of $T$.  
Using the perturbed continuity of $T$ with respect to the exact extended $b$-metric $d_{\zeta}$, 
the convergence in \eqref{limit-v} implies
\[
\lim_{n\to+\infty} d_{\zeta}(T v_n, T\vartheta)=0.
\]
Since $T v_n = v_{n+1}$, we obtain
\begin{equation}\label{limit-T}
\lim_{n\to+\infty} d_{\zeta}(v_{n+1}, T\vartheta)=0.
\end{equation}
Because $d_{\zeta}$ is a genuine extended $b$-metric on $X$, limits are unique; 
comparing \eqref{limit-v} and \eqref{limit-T} yields
\[
T\vartheta = \vartheta.
\]
Thus $\vartheta$ is indeed a fixed point of $T$.

\medskip
We next show that this fixed point is the only one.  
Assume, to the contrary, that $p,q\in X$ are two distinct fixed points of $T$.  
Applying the contractive condition to $(p,q)$ gives
\[
\mathcal{D}_{\zeta}(p,q)
= \mathcal{D}_{\zeta}(Tp,Tq)
\le c\, \mathcal{D}_{\zeta}(p,q).
\]
Writing $\mathcal{D}_{\zeta}=d_{\zeta}+\hbar$, we get
\[
d_{\zeta}(p,q)+\hbar(p,q)
\le c\,\big( d_{\zeta}(p,q)+\hbar(p,q) \big).
\]
Since $p\neq q$, the quantity $d_{\zeta}(p,q)+\hbar(p,q)$ is strictly positive, and 
the above inequality forces $c\ge 1$, contradicting the assumption $0\le c<1$.
Hence, the fixed point $\vartheta$ is unique. This completes the proof of the theorem.
\end{proof}

\end{document}
